% Generated draft: requires mathematical and source review.
\documentclass[11pt]{article}
\usepackage[margin=25mm]{geometry}
\usepackage{fontspec}
\setmainfont{Latin Modern Roman}
\usepackage{amsmath,amssymb,mathtools,unicode-math}
\setmathfont{Latin Modern Math}
\usepackage{enumitem,needspace,xurl,hyperref}
\setlength{\parindent}{0pt}
\setlength{\parskip}{5pt}
\setlength{\emergencystretch}{3em}
\newcommand{\N}{\mathbb{N}}
\newcommand{\Z}{\mathbb{Z}}
\newcommand{\Q}{\mathbb{Q}}
\newcommand{\R}{\mathbb{R}}
\newcommand{\C}{\mathbb{C}}
\newcommand{\F}{\mathbb{F}}
\newcommand{\A}{\mathbb{A}}
\newcommand{\PP}{\mathbb{P}}
\newcommand{\E}{\mathbb{E}}
\newcommand{\eps}{\varepsilon}
\newcommand{\dd}{\,\mathrm{d}}
\begin{document}
\section*{UnsolvedMath ID: 30004595}
\textbf{Source problem code:} catalogue entry
\subsection{Definitions and mathematical statement}
Let \(f=(f(x))_{x\in\R^2}\) be a real Gaussian random field: every finite vector \((f(x_1),\ldots,f(x_n))\) has a multivariate Gaussian distribution. Stationarity means that these finite-dimensional distributions are invariant under simultaneous translation of the points \(x_i\). Define the closed-level excursion set, with the boundary level included, by
\[
E_0(f)=\{x\in\R^2:f(x)\leq 0\}.
\]
Here ``connected component'' means a maximal connected subset in the Euclidean subspace topology; a component is bounded if it is contained in some Euclidean ball. The adjective ``closed-level'' specifies the non-strict inequality and does not assert that the sample excursion set is topologically closed without a continuity hypothesis.

\textbf{Source question, with unresolved scope.} For the intended planar stationary Gaussian field \(f\), in particular the random plane wave named in the source, is it true that
\[
\mathbb P\!\left(\text{every connected component of }E_0(f)\text{ is bounded}\right)=1?
\]
Equivalently, the displayed event says that for each connected component \(C\) there exists a finite radius \(R_C\) such that \(C\subseteq B(0,R_C)\); no common radius for all components is requested.

\textbf{Unresolved qualifications.} The almost-sure formulation is the natural probabilistic interpretation of the supplied question, which does not explicitly state its probability quantifier. The record does not specify the intended class of stationary Gaussian fields, centering, nondegeneracy, sample regularity, covariance assumptions, or a particular sample version. It names the random plane wave without defining its law or normalization. Thus the display records the proposed boundedness property, rather than a fully specified universal conjecture for all stationary Gaussian fields. Its interpretation also presupposes that the displayed event is measurable for the intended model.

The source additionally asks whether the phase transition at level \(0\) is continuous and calls this equivalent to the boundedness question. It supplies neither a definition of continuity of the transition nor the hypotheses establishing this equivalence or making \(0\) the critical level for the broader class. That claimed equivalence is retained as source wording requiring review, not asserted here as a theorem.
\subsection{Short English statement}
The question asks whether, almost surely, the planar set where the intended Gaussian field satisfies \(f\leq 0\) has only bounded connected components, especially for the random plane wave. The source describes this as continuity of the transition at level \(0\), but the field hypotheses and the meaning and justification of that equivalence remain unspecified.
\subsection{Sources}
\begin{itemize}
\item Local UnsolvedMath registry, supplied complete source record: permanent ID \(30004595\), source problem code \texttt{catalogue entry}, title ``Continuity of the Zero-Level Gaussian Percolation Transition.'' Collection: Oberwolfach Reports Open Problems. The supplied proposed year is 2021; no individual proposer is supplied.
\item Supplied source citation: \emph{Spatial Networks and Percolation (2022)}, Oberwolfach Reports, DOI 10.4171/owr/2021/4, pp.\ 149--224. \url{https://doi.org/10.4171/owr/2021/4}. The literature list also labels this reference ``Oberwolfach Report 4/2021.'' The relationship between the supplied 2022 title date and the 2021 report designation requires source verification.
\item Supplied literature reference: Muirhead--Rivera--Vanneuville--Koehler-Schindler, \emph{Annals of Probability}. \url{https://arxiv.org/abs/2010.11770}. The record supplies no article title or further journal bibliographic details.
\item Literature assessment dated exactly \texttt{2026-08-22}: the supplied curation labels the problem open, states that the critical level is established for the random plane wave, and says that boundedness of zero-level components remains unresolved and that the cited Gaussian-percolation work leaves absence of percolation at criticality open. This is supplied dated curation, not an independently verified literature assessment. The background expressly requests independent verification of time-sensitive status and cited claims.
\item The registry gives statement status \texttt{corrected\_verified}, qualified by a publication-format audit that preserved mathematical content while repairing notation or typography. This supplied label does not constitute independent mathematical verification of the formulation.
\end{itemize}
\end{document}
